\chapter{Resultados Esperados}

De acuerdo con el alcance definido y la ventana de ejecución experimental, este proyecto se compromete con la entrega de un sistema tecnológico funcional a escala de laboratorio. El propósito central es validar la arquitectura híbrida (Sensores + EDO + IA) no solo como herramienta de medición, sino como un instrumento activo para la **gestión sostenible**, estableciendo las bases para la reducción de la huella de carbono en implementaciones industriales.

A continuación se detallan los entregables técnicos, científicos y ambientales comprometidos:

\section{Prototipo Funcional de Monitoreo IoT}
Se entregará un \textbf{sistema de adquisición de datos operativo y calibrado}, diseñado específicamente para resistir las condiciones del entorno de bioconversión.
\begin{itemize}
    \item \textbf{Entregable:} Nodo de sensores integrado (ESP32) capaz de transmitir telemetría continua de \ce{CO2}, \ce{CH4}, temperatura y humedad, en los intervalos de medición, permitiendo la correlación de variables ambientales con la actividad metabólica.
    \item \textbf{Criterio de Éxito:} El sistema garantizará la continuidad de la transmisión de datos durante intervalos de tiempo definidos en ciclos completos (14 días), con protocolos de reconexión automática ante fallos de red y una caracterización documentada de la deriva instrumental (se aceptará un MAPE $\leq 15$\,\% respecto a la tendencia de equipos de referencia, considerando la saturación de los sensores \textit{low-cost}).
\end{itemize}

\section{Modelos Matemáticos y Algorítmicos (Arquitectura Híbrida)}
Se formulará y validará un esquema de modelado dual que combine principios biológicos con inteligencia computacional.
\begin{itemize}
    \item \textbf{Componente Mecanístico:} Un modelo EDO parametrizado que describa la curva teórica de emisiones de \ce{CO2} basada en el crecimiento larval ideal.
    \item \textbf{Componente de IA:} Un algoritmo de aprendizaje automático (evaluando regresión, árboles de decisión y redes neuronales ligeras) entrenado con datos experimentales para corregir las discrepancias del modelo teórico y detectar anomalías.
    \item \textbf{Criterio de Éxito:} Demostración estadística de que el modelo híbrido reduce el error de estimación en comparación con el uso del modelo teórico por sí solo (reducción del RMSE $\geq 15$\,\%).
\end{itemize}

\section{Plataforma de Datos y Visualización}
Implementación de una infraestructura digital básica pero escalable para la gestión de la información experimental.
\begin{itemize}
    \item \textbf{Entregable:} Despliegue de una base de datos y un tablero de control (\textit{Dashboard}) web.
    \item \textbf{Funcionalidad:} La plataforma permitirá la visualización de curvas de tendencia en tiempo real, el almacenamiento histórico seguro y la descarga de reportes estructurados para análisis posterior.
\end{itemize}

\section{Piloto de Sistema MRV (Medición, Reporte y Verificación)}
Desarrollo de un protocolo digital preliminar para la trazabilidad ambiental.
\begin{itemize}
    \item \textbf{Alcance:} El sistema generará reportes de Emisión por Lote, calculando las emisiones directas de GEI (\ce{CO2-eq}/kg residuo procesado) basadas en los datos procesados.
    \item \textbf{Limitación Declarada:} Esta versión no constituirá una certificación legal, sino una herramienta de validación interna para demostrar la viabilidad técnica del monitoreo continuo frente a los métodos de estimación estáticos.
\end{itemize}

\section{Generación de Estrategias para la Mitigación de GEI}
Como aporte directo a la sostenibilidad del proceso, se derivarán lineamientos operativos a partir del análisis de datos.
\begin{itemize}
    \item \textbf{Entregable:} Identificación de "Ventanas Operativas Limpias". Se documentarán los rangos críticos de temperatura y humedad que correlacionan con la aparición de picos de metano (\ce{CH4}) en el sustrato específico evaluado.
    \item \textbf{Impacto Ambiental:} Esta caracterización permitirá proponer recomendaciones de manejo (ej. frecuencia de volteo o aireación) orientadas a \textbf{prevenir} la anaerobiosis, transformando el sistema de una herramienta pasiva de monitoreo a un instrumento activo para la mitigación de emisiones.
\end{itemize}

\section{Delimitación del Alcance y Restricciones Técnicas}
Con el fin de garantizar la rigurosidad académica, se declaran explícitamente las siguientes fronteras operativas del proyecto:

\begin{itemize}
    \item \textbf{Precisión Instrumental:} Dado el uso de sensores de bajo costo en ambientes de alta saturación hídrica, se aceptarán márgenes de tolerancia superiores a los de equipos analíticos de laboratorio, compensados mediante algoritmos de software.
    \item \textbf{Escala de Datos:} Los modelos de Inteligencia Artificial tendrán un carácter exploratorio, entrenados con un volumen de datos limitado a los ciclos experimentales del proyecto, lo cual podría restringir su capacidad de generalización inmediata a otros sustratos no testeados.
    \item \textbf{Nivel de Automatización:} El sistema se centrará en el monitoreo y la generación de alertas (lazo abierto); no se implementará control automático de actuadores (ventiladores/riego) en esta fase.
\end{itemize}

Esta propuesta de resultados busca equilibrar la innovación tecnológica con la factibilidad operativa, entregando una solución sólida, validada experimentalmente y documentada, que cumple con los requisitos de una tesis doctoral sin exceder las capacidades logísticas disponibles.